\section{Inflation factor} \label{sec:infl}

%

The main goal of this section is to give an explicit formula for the inflation factor of the substitution $\varrho^{  }_{\bf{a}}$ provided $\mathbf a$ satisfies conditions \ref{def:a1}--\ref{def:a3}. This is done in Proposition~\ref{prop:length} below. First, we introduce an auxiliary parameter $\mu:=\mu(\mathbf a)$ which is the unique number in the interval $(0,1)$ that satisfies

\begin{equation}\label{eq:mu}
\frac{1}{\mu}=\sum_{i=0}^\infty a_i \mu^i.
\end{equation}

Existence and uniqueness follow from the observation that both parts of the equality are continuous as functions of $\mu$ on $(0,1)$ and the left-hand side is decreasing from $\infty$ to $0$ while the right-hand side is increasing from $a_0$ to $\infty$.
 
With this $\mu$, we set $$\lambda:=\mu+\frac{1}{\mu}>2.$$
Our first goal is to show that we can find an appropriate sequence $\mathbf a$ to get every possible~$\lambda$.




\begin{lem}\label{lem:inflation}
For every prescribed $\lambda>2$ we can find a suitable sequence $\mathbf a$ satisfying conditions \textnormal{\ref{def:a1}--\ref{def:a3}} such that $\lambda=\mu+\frac{1}{\mu}$ and $\mu$ satisfies Eq.~\eqref{eq:mu}.
\end{lem}

\begin{proof}
Let $0< \mu < 1$ be such that $\mu+\dfrac{1}{\mu}=\lambda$. 
First, we choose a positive integer $C$ such that $\mu+\mu^C\leq 1$. Then
$$\frac{1}{\mu}\geq \frac{1}{1-\mu^C}=\sum_{i=0}^\infty\mu^{iC}.$$
Let $\mu':=\frac{1}{\mu}-\frac{1}{1-\mu^C}\geq 0$ be the difference. We will show that it is possible to represent $\mu'$ as a sum of powers of $\mu$ with bounded nonnegative coefficients. The case of $\mu'=0$ is obvious as we can use zeros as the coefficients so we treat only the case $\mu'>0$ below.

Let $N$ be a positive integer such that $$\frac{1}{\mu}\leq \dfrac{N}{1-\mu}=\sum_{i=0}^\infty N\mu^i.$$
We consider the following family $F^0$ consisting of sequences of integers from $\left\{0,1,\ldots,N\right\}$ whose associated $\mu$-series sum is not greater than $\mu'$, i.e.,
$$F^0:=\left\{(b_i)_i\in \{0,1,\ldots,N\}^{\N_0} \colon \sum_{i=0}^\infty b_i\mu^i\leq \mu' \right\}.$$
The family $F^0$ contains infinitely many sequences. We recursively choose the sequence of subfamilies
$$F^0 \supseteq F^1 \supseteq F^2 \ldots$$ such that the $i$th term of every sequence in $F^{i+1}$ is the maximal number $c_i$ in $\{0,1,\ldots,N\}$ with the condition that $F^i$ contains infinitely many sequences with the $i$th term equal to $c_i$. 

We claim that ${\displaystyle \mu'=\sum_{i=0}^{\infty}c_i\mu^i}$. Indeed, the right-hand side cannot be greater than $\mu'$ because for every $i$, the sequence $(c_0,c_1,\ldots,c_i,0,0,\ldots)$ belongs to $F^i$.  

If ${\displaystyle \mu'>\sum_{i=0}^{\infty}c_i\mu^i}$,  we consider two cases. If there are infinitely many $c_i$s that are less than $N$, then for one sufficiently large $i$, we can swap $c_i$ with $c_i+1$ and keep the sum of the series less than $\mu'$. This contradicts the choice of $c_i$ since, in that case, there are infinitely many sequences in $F^i$ with $i$th term equal to $c_i+1 \leq N$.

In the second case, the sequence $(c_i)_i$ stabilizes at the value $N$. Let $j$ be the index such that
$$c_j<N=c_{j+1}=c_{j+2}=\ldots.$$ We claim that the sequence $(c_0,c_1,\ldots,c_{j-1},c_j+1,0,0,\ldots)$ gives the series with sum less than $\mu'$ and therefore $c_j$ was not chosen to be maximal. Indeed, by our assumption ${\displaystyle \mu'>\sum_{i=0}^{\infty}c_i\mu^i=\sum_{i=0}^{j}c_i\mu^i+\sum_{i=j+1}^{\infty}N\mu^i}$. By the choice of $N$, ${\displaystyle \mu^j\leq \sum_{i=j+1}^{\infty}N\mu^i}$. Combining these two inequalities, we get the claim and a contradiction with our assumption that $\mu'$ is greater than the sum of the series.

To finish the proof, we write
$$\frac{1}{\mu}=\sum_{i=0}^\infty\mu^{iC}+\mu'=\sum_{i=0}^\infty\mu^{iC}+\sum_{i=0}^\infty c_i\mu^i.$$
The first series ensures that properties \ref{def:a2} and \ref{def:a3} hold, and the second series guarantees property \ref{def:a1} for the sum.
\end{proof}

In the previous section, we have established that we can apply Theorem \ref{thm:primrecquasi} to our substitution.
In particular, it means there exists a natural tile length function associated with $\varrho^{ }_{\mathbf{a}}$, which is strictly positive.
Next, we show that $\lambda=r(M)$ and give explicit natural lengths for all letters in the subalphabet $\iota_{\ab}(\N_0)=\left\{[i]\colon i\in \N_0\right\}$ of isolated points. Recall that
$$\mathbf A=\left( 
\begin{array}{cccccc}
a_0& a_1+1& a_2 & a_3 & a_4 & \dots\\
1 & 0 & 1 & 0 & 0 & \ldots\\
0 & 1 & 0 & 1 & 0 & \ldots\\
0 & 0 & 1 & 0 & 1 & \ldots\\
0 & 0 & 0 & 1 & 0 & \ldots\\
\vdots & \vdots & \vdots & \vdots & \vdots & \ddots
\end{array}
\right).$$
Define $\ell:\iota_{\ab}(\N_0)\longrightarrow \R^+$ to be $\ell([0]):=1$, and for every $k>0$
\begin{equation}\label{eq:length}
\ell([k])=\mu^k+\sum_{j=1}^k\sum_{i=j}^\infty a_i\mu^{i+k+1-2j}.
\end{equation}


\begin{lem}\label{lem:unifcon}
The function $\ell\/\colon\/\iota_{\ab}(\N_0)\longrightarrow \R^+$ in Eq.~\eqref{eq:length} is uniformly continuous on\/ $\iota_{\ab}(\N_0)$ with respect to the topology on\/ $\mathcal{S}$.
\end{lem}

\begin{proof}
Let $[k],[k+t]\in \iota_{\ab}(\N_0)$ be such that $t>0$ and 
$d([k],[k+t])<\frac{1}{2^{n}}$ for some $n\in\N$. We show that $|\ell([k])-\ell([k+t])|<\varepsilon(n)$. First, note that $\sigma^{k}(\mathbf{0})$ and $\sigma^{k+t}(\mathbf{0})$ agree on $[-n,n]$ as the distance between these points in $\mathcal S$ is at most $\frac{1}{2^{n+1}}$. This means 
\begin{equation}\label{eq:n-close}
a_{k+r}=a_{k+t+r} \quad\text{for}\quad -n\leqslant r\leqslant n.
\end{equation} 
Consequently $n\leqslant k$ as $\sigma^k(\mathbf{0})$ has $*$ at the position $-k-1$, but all further iterations of $\sigma$ have non-$*$ entries there.


We can then split the double sum on the right hand-side of Eq.~\eqref{eq:length} into three parts:
\[
\ell([k])=\mu^k+\underbrace{\sum_{j=1}^{k-n}\sum_{i=j}^\infty a_i\mu^{i+k+1-2j}}_{\textnormal{(I)}}+\underbrace{\sum_{j=k-n+1}^k\sum_{i=k+n+1}^\infty a_i\mu^{i+k+1-2j}}_{\textnormal{(II)}}+\underbrace{\sum_{j=k-n+1}^k\sum_{i=j}^{k+n} a_i\mu^{i+k+1-2j}}_{\textnormal{(III)}}.
\]
Note that the component (III) is the same for both $\ell([k])$ and $\ell([k+t])$ because of Eq.~\eqref{eq:n-close}, and hence this component vanishes in $|\ell([k])-\ell([k+t])|$.
We now determine $n$-dependent bounds for (I). 
\begin{align*}
\sum_{j=1}^{k-n}\sum_{i=j}^\infty a_i\mu^{i+k+1-2j}& \leqslant N \sum_{j=1}^{k-n}\sum_{i=j}^\infty \mu^{i+k+1-2j}= N \sum_{j=1}^{k-n}\sum_{i=0}^\infty \mu^{i+k+1-j}\\
&=N \sum_{j=1}^{k-n} \frac{\mu^{k+1-j}}{1-\mu}=\frac{N}{(1-\mu)}(\mu^{n+1}+\cdots+ \mu^k)\leqslant \frac{N\mu^n}{(1-\mu)^2}.
\end{align*}
Here, the first inequality follows from \ref{def:a1}. One can carry out an analogous calculation for (II), which yields
\[
\sum_{j=k-n+1}^{k}\sum_{i=k+n+1}^\infty a_i\mu^{i+k+1-2j}\leqslant \frac{N\mu^{n+2}}{(1-\mu)^2}.
\]
Combining these estimates with the fact that $\mu^{k}-\mu^{k+t}< \mu^n$ gives us
\[
\left|\ell([k])-\ell([k+t])\right|<\left(1+\frac{2N}{(1-\mu)^2}+\frac{2\mu^2 N}{(1-\mu)^2}\right)\mu^n,
\]
from which the claim is immediate as the right-hand side goes to 0 as $n$ grows.
\end{proof}



Recall that $\mathcal{K}$ is the positive cone of $C(\mathcal{A})$. 



\begin{prop}\label{prop:length}
The function $\ell$ in Eq.~\eqref{eq:length}
is a left eigenvector of $\mathbf{A}$ corresponding to $\lambda$. Moreover, it extends to a strictly positive, continuous natural length function on $\mathcal{A}$ and $\lambda=r(M)$ is the associated inflation factor. 
\end{prop}

\begin{proof}
We first show that the (infinite) vector $\ell_\mathcal A:=(\ell([0]),\ell([1]),\ell([2]),\ldots)$ of natural tile lengths is a left eigenvector of the matrix $\mathbf A$ corresponding to eigenvalue $\lambda=\mu+\frac{1}{\mu}$, so $\ell_\mathcal A\mathbf A=\lambda \ell_\mathcal A$.

We consider the product in the left-hand side column by column. Particularly, the first column yields 
$a_0\ell([0])+\ell([1])=\lambda \ell([0])$. Since $\ell([0])=1$, one has
$$\ell([1])=\lambda -a_0=\mu+\frac{1}{\mu}-a_0=\mu+\sum_{i=0}^\infty a_i\mu^i-a_0=\mu+\sum_{i=1}^\infty a_i\mu^i,$$
which satisfies the statement of this proposition.

After that, looking at the $k$th column for $k>1$ we get
$$a_{k-1}\ell([0])+\ell([k-2])+\ell([k])=\lambda \ell([k-1]),$$
or
$$\ell([k])=\lambda \ell([k-1])-\ell([k-2])-a_{k-1}.$$

In particular, for $k=2$,
\begin{multline*}\ell([2])=\lambda \ell([1])-\ell([0])-a_1=\left(\mu+\frac{1}{\mu}\right)\left( \mu+\sum_{i=1}^\infty a_i\mu^i\right)-1-a_1=\\
=\mu^2+1+\sum_{i=1}^\infty a_i\mu^{i+1}+\sum_{i=1}^\infty a_i\mu^{i-1}-1-a_1=\\=
\mu^2+\sum_{i=1}^\infty a_i\mu^{i+1}+\sum_{i=2}^\infty a_i\mu^{i-1}=\mu^2+\sum_{j=1}^2\sum_{i=j}^\infty a_i\mu^{i+3-2j}.
\end{multline*}

For all other $k$, the computations can be carried out by induction. Using $\lambda=\mu+\frac{1}{\mu}$ and expanding the product $\lambda\ell([k-1])$ we get
\begin{multline*}\ell([k])=\lambda \ell([k-1])-\ell([k-2])-a_{k-1}=\\=
\left(\mu+\frac{1}{\mu}\right)\left(\mu^{k-1}+\sum_{j=1}^{k-1}\sum_{i=j}^\infty a_i\mu^{i+k-2j}\right)-\left(\mu^{k-2}+\sum_{j=1}^{k-2}\sum_{i=j}^\infty a_i\mu^{i+k-1-2j}\right)-a_{k-1}=\\=
\mu^k+\mu^{k-2}+\sum_{j=1}^{k-1}\sum_{i=j}^\infty a_i\mu^{i+k+1-2j}+\sum_{j=1}^{k-1}\sum_{i=j}^\infty a_i\mu^{i+k-1-2j}-\mu^{k-2}-\sum_{j=1}^{k-2}\sum_{i=j}^\infty a_i\mu^{i+k-1-2j}-a_{k-1}.
\end{multline*}

Here, we can cancel $\mu^{k-2}$ and also cancel the double sum with negative sign leaving only the term with $j=k-1$ from the second double sum. Simplifying, the expression, we get
$$=
\mu^k+\sum_{j=1}^{k-1}\sum_{i=j}^\infty a_i\mu^{i+k+1-2j}+\sum_{i=k-1}^\infty a_i\mu^{i+k-1-2(k-1)}-a_{k-1}.$$
Note that the first term of the second sum is $a_{k-1}$ so we can cancel that as well. Additionally, the power of $\mu$ in the second sum is $i+k-1-2(k-1)=i+k+1-2k$. These adjustments give us the final expression
$$=\mu^k+\sum_{j=1}^{k-1}\sum_{i=j}^\infty a_i\mu^{i+k+1-2j}+\sum_{i=k}^\infty a_i\mu^{i+k+1-2k}=\mu^k+\sum_{j=1}^k\sum_{i=j}^\infty a_i\mu^{i+k+1-2j}$$
as needed.

We now show that $\lambda=r(M)$. Since $\rho^{ }_{\mathbf a}$ is primitive and
$M/r(M)$ is quasicompact, the spectral radius $r(M)$ of the substitution operator $M$ is an eigenvalue with a strictly positive, continuous eigenfunction $\ell$;  see \cite[Thm. 5.4(B)]{MRW2} and  \cite[Thm.~V.5.2 and Prop.~V.5.6]{Sch}. Moreover, $r(M)$ is the only eigenvalue of $M$ with an eigenfunction in $\mathcal{K}$; see \cite[Thm.~V.5.2(iv)]{Sch}.  

Since the subalphabet $\iota_{\mathbf{a}}(\N_0)$ is dense in $\mathcal{A}$ and since $\ell$ is uniformly continuous  on $\iota_{\mathbf{a}}(\N_0)$ by Lemma~\ref{lem:unifcon}, it follows that $\ell$ extends to a unique continuous function on $\mathcal{A}$, which must necessarily be in the positive cone $\mathcal{K}$, and must be an eigenfunction of $M$ corresponding to the eigenvalue $\lambda$. It then follows that $\lambda=\mu+\frac{1}{\mu}$ must be the spectral radius $r(M)$ of the substitution operator. 
\end{proof}


\begin{rem}
The second claim in the previous result is an infinite-alphabet analogue of the Perron eigenvalue condition by Lagarias and Wang for inflation functional equations admitting weak Delone set solutions in $\mathbb{R}^{d}$, which states $\det(A)=\lambda(S)$ where $\lambda(S)$ is the Perron--Frobenius eigenvalue of the subdivision matrix $S$ and $A$ is the inflation map; see~\cite{LW}. 
\end{rem}


Our next goal is to prove Theorem \ref{thm:allfactors}. We start by giving the formal definition of a Delone set, which is needed in the theorem. A point set $\varLambda$ in $\mathbb{R}^d$ is called a Delone set if it is \begin{itemize}
\item uniformly discrete, i.e., there is a positive radius $r$ such that the balls of radius $r$ centered at points of $\varLambda$ form a packing in $\R^d$, and  
\item relatively dense, i.e., there is a positive radius $R$ such that the balls of radius $R$ centered at points of $\varLambda$ form a covering of $\R^d$.
\end{itemize}


\begin{proof}[Proof of Theorem~\textnormal{\ref{thm:allfactors}}]
For a given $\lambda$, we can find an appropriate sequence using the results of Lemma \ref{lem:inflation}. The resulting alphabet is infinite and compact. It remains to show that every element of the subshift gives rise to a Delone set. This follows automatically from the continuity and strict positivity of $\ell$, which means there exist constants $C_1,C_2>0$ such that for every $a\in \mathcal{A}$, $C_1\leq \ell(a)\leq C_2$ where $\ell$ is the natural length function. The lower bound implies that the point set is uniformly discrete, while the upper bound guarantees that it is relatively dense. 

%


In fact, we get an explicit lower bound from  property \ref{def:a3}, which states that for every $k>0$, at least one of the coefficients $a_k,\ldots,a_{k+C}$ is not zero. This means the internal sum in Eq.~\eqref{eq:length} for $j=k$ is at least 
$$a_k\mu+a_{k+1}\mu^2+\ldots+a_{k+C}\mu^{C+1}.$$
Thus for every letter $[k], k\in\N_0$, $\ell([k])\geq \mu^{C+1}$. 
This completes the proof.
\end{proof}
